\documentclass{article}
\usepackage[UTF8]{ctex}
\usepackage{amsmath,amssymb,amsfonts}
\usepackage{enumitem}
\usepackage{geometry}
\geometry{a4paper, margin=1in}
% === 错题本排版优化（不参与题目标识） ===
\setlength{\parindent}{0pt}
\setlength{\parskip}{0.35em plus 0.10em minus 0.08em}
\setlength{\emergencystretch}{3em}
\setlist{itemsep=0.25em, topsep=0.40em, parsep=0pt}
\allowdisplaybreaks[2]
% === 错题本排版优化结束 ===
\begin{document}
\section*{天津大学2025年数学分析真题}
\section*{59．天津大学2025年数学分析试卷（回忆版）}

\subsection{题目1}
1．证明： \(\int_{1}^{+\infty} \frac{\sin \pi x}{x} \mathrm{~d} x<0\) ．

\subsection{题目2}
2．\(f(x)\) 是定义在 \([a, b]\) 上的函数，且每一个点处极限都存在．证明：\(f(x)\) 在 \([a, b]\) 上有界．

\subsection{题目3}
3．设 \(f(x)\) 在 \([0, \pi]\) 上有连续可微，且 \(\int_{0}^{\pi} f(x) \mathrm{d} x=0\) ．证明：

\[
\int_{0}^{\pi} f^{2}(x) \mathrm{d} x \leq \int_{0}^{\pi}\left[f^{\prime}(x)\right]^{2} \mathrm{~d} x .
\]

\subsection{题目4}
4．已知 \(f(x)\) 在 \([0,1]\) 上二阶可导，且 \(f(0)=0, f(1)=3, \min _{x \in[0,1]} f(x)=-1\) ．证明：存在 \(c \in(0,1)\) ，使得

\[
f^{\prime \prime}(c) \geq 18
\]

\subsection{题目5}
5．设 \(x, y>0\) ，证明：\(\frac{1}{4}\left(x^{2}+y^{2}\right) \leq e^{x+2 y-2}\) ．

\subsection{题目6}
6．设 \(a_{n}, b_{n} \rightarrow 0(n \rightarrow \infty)\) ，且 \(b_{1}>b_{2}>b_{3}>\cdots\) ，若 \(\lim _{n \rightarrow \infty} \frac{a_{n}-a_{n-1}}{b_{n}-b_{n-1}}=A\) ，证明： \(\lim _{n \rightarrow \infty} \frac{a_{n}}{b_{n}}=A\) ．

\subsection{题目7}
7．设 \(f(x)\) 是 \([0,+\infty)\) 上的可导函数，且 \(\lim _{x \rightarrow+\infty} f^{\prime}(x)=A(A\) 为有限数或 \(+\infty)\) ，证明：\(f(x)\) 在\([0,+\infty)\) 上一致连续的充要条件是 \(A\) 为有限数。

\subsection{题目8}
8．设 \(\mathrm{x} \in \mathbb{R}^{3}\) ，定义范数 \(N(\mathbf{x}): \mathbb{R}^{3} \rightarrow \mathbb{R}\) 若满足

\begin{itemize}
  \item[（1）] \(N(\mathbf{x}) \geq 0\) ，且仅当 \(\mathbf{x}=\mathbf{0}\) 时，有 \(N(\mathbf{x})=0\) 。
  \item[（2）] \(N(\mathbf{x}+\mathbf{y}) \leq N(\mathbf{x})+N(\mathbf{y}), \forall \mathbf{x}, \mathbf{y} \in \mathbb{R}^{3}\) ．
  \item[（3）] \(N(k \mathbf{x})=|k| N(\mathbf{x}), \forall \mathbf{x} \in \mathbb{R}^{3}, k \in \mathbb{R}\) ．
\end{itemize}

上述空间被称为欧氏空间。内积定义的范数 \(\|\mathrm{x}\|=\sqrt{\mathbf{x} \cdot \mathbf{x}}\) 。

\begin{itemize}
  \item[（1）] 证明：\(N(\mathbf{x})\) 在 \(B=\left\{\mathbf{x} \in \mathbb{R}^{3} \mid\|\mathbf{x}\| \leq 1\right\}\) 上有界。
  \item[（2）] \(N(\mathrm{x})\) 在 \(\mathbb{R}^{3}\) 上连续．
  \item[（3）] 对任意的 \(\mathrm{x} \in \mathbb{R}^{3}\) ，记 \(r=\|\mathrm{x}\|\) ，存在 \(a, b \in \mathbb{R}\) ，使得
\[
a r \leq N(x) \leq b r
\]
\end{itemize}

\subsection{题目9}
9．记 \(H\) 为黑塞矩阵。\\
（1）在点 \((a, b)\) 处，已知 \(f_{x}(a, b)=f_{y}(a, b)=0, H=\left(\begin{array}{ll}f_{x x}(x, y) & f_{x y}(x, y) \\ f_{y x}(x, y) & f_{y y}(x, y)\end{array}\right)\) ，且 \(H\) 在 \((a, b)\) 处正定，证明：\((a, b)\) 为 \(f\) 的严格极小值点。\\
（2）若 \(H\) 在 \(\mathbb{R}^{2}\) 上处处正定，则 \(f\) 在 \(\mathbb{R}^{2}\) 上至多有一个极值点。

\subsection{题目10}
10．设 \(B\) 为单位球，\(f\) 在 \(B\) 上有定义，\(v=\left(v_{1}, v_{2}, v_{3}\right), \nabla \cdot v=0\) ．在 \(B\) 上有 \(v=0\) ．

\begin{itemize}
  \item[（1）] 求证： \(\int_{B} \nabla f \cdot v \mathrm{~d} x=0\) 。
  \item[（2）] 求证： \(\int_{B} v_{1} \mathrm{~d} x=0\) ．
\end{itemize}
\end{document}
